\chapter*{Abstract}
\addcontentsline{toc}{chapter}{Abstract}

Academic institutions accumulate large volumes of student performance data, yet these records are often used only for administrative reporting rather than timely academic intervention. Manual interpretation of grade sheets, backlog status, attendance compliance, and semester-wise trends is time-consuming for both students and faculty advisors. As a result, many at-risk students are identified only after academic decline becomes severe.

This report presents \textbf{SARS}, an Explainable Agentic Retrieval-Augmented Generation based Student Academic Performance Risk Assessment and Advisory System. The proposed system combines document-based grade ingestion, structured academic data storage, explainable risk scoring, and evidence-grounded academic advisory support in a single platform. Official grade sheets are uploaded in image or PDF form and processed through a vision-based extraction workflow, after which the extracted data are verified by the user before being stored. A transparent SARS score is then computed using weighted risk factors derived from CGPA, backlog history, attendance compliance, and semester trend behavior.

The system further includes a Retrieval-Augmented Generation pipeline that converts student records into typed knowledge chunks, stores them in a vector database, and retrieves relevant evidence during advisory conversations. This enables the advisory component to provide personalized recommendations grounded in the student's own academic record rather than generic responses. Student-facing dashboards support marksheet upload, academic record review, risk visualization, attendance tracking, and academic advisory. Teacher-facing functionality supports monitoring of multiple students, drill-down inspection, and intervention logging.

The project demonstrates how explainability, institutional policy alignment, and conversational AI can be combined to support early academic intervention. The report presents the background, architecture, methodology, implementation, results, and future scope of the SARS system in a thesis-style format.

\clearpage
